% Part: first-order-logic
% Chapter: syntax-and-semantics
% Section: formation-sequences

\documentclass[../../../include/open-logic-section]{subfiles}

\begin{document}

\olfileid{pl}{syn}{fseq}

\olsection{నిర్మాణ క్రమాలు}

ఆగమనాత్మక నిర్వచనం ద్వారా !!{formula}sను నిర్వచించడం, దానికి
పూరకంగా ఆగమనం ద్వారా !!{formula}s ధర్మాలను నిరూపించడం సొగసైన,
సమర్థమైన విధానం. అయితే !!{formula}s నిర్మాణాన్ని కింది నుంచి పైకి,
మెట్టుమెట్టుగా పరిశీలించడం కూడా ఉపయోగకరం కావచ్చు. ఇక్కడ
\emph{నిర్మాణ క్రమం} అనే భావనతో అదే చేస్తాం.

\begin{defn}[సూత్రాల నిర్మాణ క్రమాలు]
\ollabel{defn:fseq-frm}
భాష~$\Lang L_0$లోని సంకేతాలతో ఏర్పడిన సంకేతమాలల పరిమిత క్రమం
$\tuple{!A_0,\dotsc,!A_n}$, $!A$కు \emph{నిర్మాణ క్రమం} కావడం
అంటే $!A \ident !A_n$, అలాగే ప్రతి $i \leq n$కు $!A_i$ ఒక పరమాణు
సూత్రమో, లేక కింది వాటిలో ఒకటి సత్యమయ్యేలా $j,k < i$లు ఉంటాయో:
\begin{enumerate}
    \tagitem{prvNot}{$!A_i \ident \lnot !A_j$.}{}%
    \tagitem{prvAnd}{$!A_i \ident (!A_j \land !A_k)$.}{}%
    \tagitem{prvOr}{$!A_i \ident (!A_j \lor !A_k)$.}{}%
    \tagitem{prvIf}{$!A_i \ident (!A_j \lif !A_k)$.}{}%
    \tagitem{prvIff}{$!A_i \ident (!A_j \liff !A_k)$.}{}%
\end{enumerate}
\end{defn}

\begin{ex}
\[
\tuple{
    \Obj p_0,
    \Obj p_1,
    (\Obj p_1 \land \Obj p_0),
    \lnot (\Obj p_1 \land \Obj p_0)
}
\]
అనేది $\lnot (\Obj p_1 \land \Obj p_0)$ యొక్క ఒక నిర్మాణ క్రమం;
కిందిదీ అలాంటి క్రమమే:
\[
\tuple{
    \Obj p_0,
    \Obj p_1,
    \Obj p_0,
    (\Obj p_1 \land \Obj p_0),
    (\Obj p_0 \lif \Obj p_1),
    \lnot (\Obj p_1 \land \Obj p_0)
}.
\]
%
రెండవ ఉదాహరణలో కనిపించినట్లుగా, నిర్మాణ క్రమాల్లో “వ్యర్థ” భాగాలు
ఉండవచ్చు: అవి అనవసరమైన లేదా నిర్మాణానికి తోడ్పడని సూత్రాలు.
\end{ex}

\begin{prop}
\ollabel{prop:formed}
$\Frm[L_0]$లోని ప్రతి !!{formula}~$!A$కు ఒక నిర్మాణ క్రమం ఉంటుంది.
\end{prop}

\begin{proof}
$!A$ పరమాణు సూత్రం అనుకుందాం. అప్పుడు క్రమం~$\tuple{!A}$,
$!A$కు నిర్మాణ క్రమం.
%
ఇప్పుడు $!B$, $!C$లకు వరుసగా
$\tuple{!B_0,\dotsc,!B_n}$, $\tuple{!C_0,\dotsc,!C_m}$ అనే నిర్మాణ
క్రమాలు ఉన్నాయని అనుకుందాం.
%
\begin{enumerate}
  \tagitem{prvNot}{$!A \ident \lnot !B$ అయితే,
      $\tuple{!B_0,\dotsc,!B_n,\lnot !B_n}$ అనేది $!A$కు నిర్మాణ
      క్రమం.}{}
  \tagitem{prvAnd}{$!A \ident (!B \land !C)$ అయితే,
      $\tuple{!B_0,\dotsc,!B_n,!C_0,\dotsc,!C_m,(!B_n \land !C_m)}$
      అనేది $!A$కు నిర్మాణ క్రమం.}{}
  \tagitem{prvOr}{$!A \ident (!B \lor !C)$ అయితే,
      $\tuple{!B_0,\dotsc,!B_n,!C_0,\dotsc,!C_m,(!B_n \lor !C_m)}$
      అనేది $!A$కు నిర్మాణ క్రమం.}{}
  \tagitem{prvIf}{$!A \ident (!B \lif !C)$ అయితే,
      $\tuple{!B_0,\dotsc,!B_n,!C_0,\dotsc,!C_m,(!B_n \lif !C_m)}$
      అనేది $!A$కు నిర్మాణ క్రమం.}{}
  \tagitem{prvIff}{$!A \ident (!B \liff !C)$ అయితే,
      $\tuple{!B_0,\dotsc,!B_n,!C_0,\dotsc,!C_m,(!B_n \liff !C_m)}$
      అనేది $!A$కు నిర్మాణ క్రమం.}{}
\end{enumerate}
!!{formula}sపై ఆగమన సూత్రం ప్రకారం ప్రతి !!{formula}కు ఒక నిర్మాణ
క్రమం ఉంటుంది.
\end{proof}

దీని తిరుగుదిశను కూడా నిరూపించవచ్చు. సూత్రాలకు మనం ఇచ్చిన రెండు
నిర్వచన పద్ధతులూ తుల్యమైనవని, అంటే అవి ఒకే ఫలితాలను ఇస్తాయని ఇది
చూపుతుంది కనుక ముఖ్యం. నిర్మాణ క్రమాల పొడవుపై సాధారణ ఆగమనాన్ని
ఉపయోగించి సూత్రాల గురించిన సిద్ధాంతాలను నిరూపించవచ్చని కూడా దీని
అర్థం.

\begin{lem}
\ollabel{lem:fseq-init}
$\tuple{!A_0,\dotsc,!A_n}$ అనేది $!A_n$కు నిర్మాణ క్రమమనీ,
$k \leq n$ అనీ అనుకుందాం. అప్పుడు
$\tuple{!A_0,\dotsc,!A_k}$ అనేది $!A_k$కు నిర్మాణ క్రమం.
\end{lem}

\begin{proof}
సాధన.
\end{proof}

\begin{thm}
\ollabel{thm:fseq-frm-equiv}
$\Frm[L_0]$ అనేది నిర్మాణ క్రమం గల భాష~$\Lang L_0$లోని అన్ని
సంకేతమాలల సమితి.
\end{thm}

\begin{proof}
భాష~$\Lang L_0$లో నిర్మాణ క్రమం గల అన్ని సంకేతమాలల సమితిని $F$
అనుకుందాం. \olref[pl][syn][fseq]{prop:formed}లో
$\Frm[L_0] \subseteq F$ అని చూశాం; ఇప్పుడు తిరుగుదిశను నిరూపిస్తాం.

$!A$కు $\tuple{!A_0,\dotsc,!A_n}$ అనే నిర్మాణ క్రమం ఉందనుకుందాం.
$n$పై బలమైన ఆగమనం ద్వారా $!A \in \Frm[L_0]$ అని నిరూపిస్తాం.
$m < n$ పొడవు గల నిర్మాణ క్రమం ఉన్న ప్రతి సంకేతమాల
$\Frm[L_0]$లో ఉంటుందనేది మన ఆగమన పరికల్పన. నిర్మాణ క్రమం నిర్వచనం
ప్రకారం, $!A_n$ పరమాణు సూత్రమో, లేక కింది వాటిలో ఒకటి సత్యమయ్యేలా
$j,k < n$లు ఉంటాయో:
\begin{enumerate}
    \tagitem{prvNot}{$!A_n \ident \lnot !A_j$.}{}%
    \tagitem{prvAnd}{$!A_n \ident (!A_j \land !A_k)$.}{}%
    \tagitem{prvOr}{$!A_n \ident (!A_j \lor !A_k)$.}{}%
    \tagitem{prvIf}{$!A_n \ident (!A_j \lif !A_k)$.}{}%
    \tagitem{prvIff}{$!A_n \ident (!A_j \liff !A_k)$.}{}%
\end{enumerate}
ఇప్పుడు సందర్భాలవారీగా తర్కిద్దాం. $!A_n$ పరమాణు సూత్రమైతే
$!A_n \in \Frm[L_0]$. బదులుగా $!A \ident
(!A_j \land !A_k)$ అనుకుందాం.
\sourcecorrection{OLTEPLSYN-003}{ఆంగ్ల మూలం ఈ వాక్యనిర్మాణ సందర్భంలో
అర్థపర తుల్యత సంకేతం $\equiv$ను వాడింది. ముందు నిర్వచించిన నిర్మాణ
క్రమాలు సంకేతమాల తాదాత్మ్యం $\ident$ను కోరుతాయి కనుక అదే సంకేతాన్ని
ఇక్కడ నిలకడగా వాడాం.}
\olref[pl][syn][fseq]{lem:fseq-init} ప్రకారం,
$\tuple{!A_0,\dotsc,!A_j}$, $\tuple{!A_0,\dotsc,!A_k}$లు వరుసగా
$!A_j$, $!A_k$లకు నిర్మాణ క్రమాలు. ఇవి $!A$ నిర్మాణ క్రమానికి నిజ
ఆద్య ఉపక్రమాలు కనుక రెండింటి పొడవూ $n$కంటే తక్కువ. అందువల్ల ఆగమన
పరికల్పన ప్రకారం $!A_j$, $!A_k$లు $\Frm[L_0]$లో ఉంటాయి; కాబట్టి
!!a{formula} నిర్వచనం ప్రకారం $(!A_j \land !A_k)$ కూడా అందులోనే
ఉంటుంది. మిగతా సందర్భాలు సమాంతర తర్కంతో వస్తాయి.
\end{proof}

\end{document}
